%% kapa: exam=1 round=24 year=2013 period=1 form=A subject=civil_law subject_ko=민법 seq=8 number=8
%% source: Q-Net 감정평가사 기출문제 (cst00309) artlSeq=5073802 file=extracted/1008/5073802/2013년 제24회 감정평가사 1차시험 기출문제/1교시(A형).hwp
%% conversion: kordoc 4.24.51 -> tools/gen1.py (auto); review=proofread-sonnet; flags=none
\begin{question}{8}
의사표시의 효력발생에 관한 설명으로 옳지 않은 것은? (다툼이 있으면 판례에 의함)
\begin{choices}
  \item 보통우편으로 발송된 우편물은 상당기간 내에 도달한 것으로 추정된다.
  \item 격지자간의 계약은 승낙의 통지를 발송한 때에 성립한다.
  \item 표의자가 의사표시를 발송한 후 사망하더라도 그 효력에 영향이 없다.
  \item 의사표시가 도달된 후에는 표의자의 상대방이 이를 요지(了知)하기 전이라도 표의자는 이를 철회하지 못한다.
  \item 의사표시의 효력발생을 위해 표의자의 상대방이 그 의사표시의 내용을 알 것까지 요구되지는 않는다.
\end{choices}
\end{question}
\answer{1}
